\newpage
\phantomsection
\label{prf:z-mod-four-has-zero-divisors}
\LRAProofFor{thm:z-mod-four-has-zero-divisors}

\begin{remark*}[Return]
\hyperref[thm:z-mod-four-has-zero-divisors]{Return to Theorem}
\end{remark*}

\begin{theorem*}[\(\mathbb{Z}/4\mathbb{Z}\) Has Zero Divisors]
In \(\mathbb{Z}/4\mathbb{Z}\), the nonzero class \([2]_4\) satisfies
\([2]_4\cdot[2]_4=[0]_4\). Therefore \(\mathbb{Z}/4\mathbb{Z}\) is not a
field.
\end{theorem*}

\begin{proof}[Professional Standard Proof]
\LRAProofBodyStart
\textbf{Professional Standard Proof.}
The class \([2]_4\) is nonzero because \(4\nmid 2\). But
\[
[2]_4\cdot[2]_4=[4]_4=[0]_4.
\]
Thus \([2]_4\) is a nonzero zero divisor. A field has no nonzero zero
divisors, so \(\mathbb{Z}/4\mathbb{Z}\) is not a field.
\end{proof}

\begin{proof}[Detailed Learning Proof]
\LRAProofBodyStart
\textbf{Detailed Learning Proof.}
In the quotient modulo \(4\), the zero class consists of integers divisible
by \(4\). Since \(2\) is not divisible by \(4\), \([2]_4\neq[0]_4\). However,
modular multiplication gives
\[
[2]_4\cdot[2]_4=[2\cdot2]_4=[4]_4=[0]_4.
\]
So a nonzero element multiplies by a nonzero element to give zero. In a field,
every nonzero element has a multiplicative inverse, which rules out such zero
divisors. Therefore \(\mathbb{Z}/4\mathbb{Z}\) is not a field.
\end{proof}

\begin{remark*}[Proof structure]
The example isolates why finite modular systems need not be fields.
\end{remark*}

\begin{dependencies}
\begin{itemize}
  \item \hyperref[def:finite-modular-number-system]{Finite Modular Number System}
  \item \hyperref[def:multiplication-modulo-n]{Multiplication Modulo \(n\)}
  \item \hyperref[def:two-sided-inverse]{Two-Sided Inverse}
\end{itemize}
\end{dependencies}

\clearpage
